\section{Cycle subgroups and explicit representatives}\label{sec:cycle-domains}

\subsection{Algebraically trivial cycles}

Write $\CH^r(X)_{\Q,\alg}$ for the rational Chow classes whose
base change to an algebraic closure is algebraically equivalent
to zero.

\begin{proposition}\label{prop:algebraic-triviality}
On every smooth projective $X/K$ and in every codimension,
\[
 \CH^r(X)_{\Q,\alg}\subseteq\cK_K^r(X).
\]
Over a number field, the companion normalized height is therefore
defined on
\[
 \CH^r(X)_{\Q,\alg}\times
 \CH^{n+1-r}(X)_{\Q,\alg},
\]
and more generally on either algebraically trivial factor paired
with the admitted complementary group $\cH$.
\end{proposition}

\begin{proof}
Algebraic equivalence is generated by connected smooth projective
curve families. After a finite extension, an algebraically trivial
class is a finite sum of
\[
 \Gamma_*([a]-[b]),\qquad
 \Gamma\in\CH^r(C\times X)_\Q,
\]
with $a,b$ rational over that finite extension on the smooth
projective parameter curve. Enlarge the extension to retain the
finite rational-equivalence witnesses as well.
The degree-zero divisor $[a]-[b]$ is in its endpoint kernel.
Kernel covariance and full-degree finite descent prove the local
inclusion. The constant-field interpretation of algebraic
triviality is equivalently obtained by norming the finite family
certificate. Apply this inclusion at every completion and then
the companion height theorem.
\end{proof}

Together with the endpoint results, this also gives the full
divisor--zero-cycle height on an arbitrary $n$-fold. The complete
one-cycle group lies in $\cH^{n-1}(X)$ and pairs with
$\cH^2(X)$. Whenever one of the geometric or model criteria above
establishes the entire codimension-two group, this becomes the
full complementary pairing.

\subsection{Finite obstruction certificates}

The finite-detector theorem of \cite{InternalDetectors} provides,
for a fixed admissible model calculation, finitely many rational
period coordinates for the actual obstruction. Their vanishing
is a certificate that a specified homological Chow class belongs
to $\cK_K^r(X)$. The normalized tests use
\[
 \eta_j=\frac{1}{[L:K]c_L}\Tr^K_{X_L/X}(b_j),\qquad
 c_L=-\log|\varpi_L|=c_K/e(L/K).
\]
Thus all model families, isotypic sectors and individually certified
classes may be combined inside the same admitted group. For a
number-field class one checks this membership at its finitely many
relevant bad places; the good places are handled by the good-model
construction. The output Green identity still has the complete
actual test space in \eqref{eq:strong-G}.

\subsection{Exact-support cancellation for complete intersections}

Let $\bar L$ be a line bundle with one fixed piecewise-smooth or
smooth metric allowed by the complete-intersection construction
on a smooth projective $X/K$, and let
$s_0,s_1$ be regular rank-$r$ tuples of sections of $L$, with
$n\geq r\geq3$. Write $Y_i=V(s_i)$ with its scheme-theoretic
cycle multiplicities, and put
\[
 Z=[Y_1]-[Y_0],\quad
 S=|Y_0|^{\an}\cup|Y_1|^{\an},\quad
 U=X^{\an}\setminus|Z|^{\an},\quad W=X^{\an}\setminus S.
\]
The complete-intersection kernel of the tropical companion
\cite{InternalTropical}, for a common allowed metric on $L$, is
\begin{equation}\label{eq:ci-kernel}
 \phi_s=\min_i(-\log\|s_i\|),\quad
 R=c_1(\bar L),\quad A_s=R-D\phi_s,\quad
 G_s=\phi_s\sum_{j=0}^{r-1}A_s^jR^{r-1-j}.
\end{equation}
It satisfies $D[G_s]+\delta_{Y_s}=[R^r]$ against all actual
tests. The raw difference is defined on $W$; cancelled components
lie in $C_U=S\cap U$.

\begin{proposition}[Supplied cancellation certificates]\label{prop:ci-cancellation}
Each of the following hypotheses supplies an explicit actual
Green form for $Z$ on $U$, satisfying \eqref{eq:strong-G}.
\begin{enumerate}[label=\textup{(\roman*)}]
\item On an ordinary analytic neighborhood of $C_U$ in $U$,
$s_1=Ms_0$ exactly, where $M$ is an invertible three-row block
plus the identity in the other rows. Two specified entries $a,d$
in its first column have nonvanishing opens covering $C_U$.
\item On a Zariski neighborhood of $C_U$, the two regular
sequences define the same ideal. Their comparison matrix on the
reduced conormal coefficient ring has a supplied finite ordered
factorization into elementary shears, one-entry units, invertible
two-row blocks, and invertible three-row blocks with the two-pivot
cover in \textup{(i)}. All row labels, coefficient data and inverses
are part of the factorization. At each fixed word position use
the same cap offsets and inner and outer partitions for every
local coefficient lift.
\end{enumerate}
\end{proposition}

\begin{proof}
For (i), the relative three-row face identities in
\cite{InternalTropical} give, for an ordinary neighborhood $N$
of $C_U$ in $U$, the identity on $N\cap W$
\[
 G_{s_1}-G_{s_0}=d'\alpha+d''\beta+E,
\]
where $\alpha,\beta$ are actual on $N\cap W$ and $E$ is
actual on $N$. Their explicit
two-pivot partition uses $v(d)-v(a)$ and includes the differentiated
partition and mixed-derivative terms. Choose the companion
relative-support cutoff $\zeta$ on $U$, with support closed in
$U$ inside this neighborhood and value one near $C_U$.
On $W$ the corrected form
\[
 G_{s_1}-G_{s_0}
 -d'\operatorname{ext}_0(\zeta\alpha)
 -d''\operatorname{ext}_0(\zeta\beta)
\]
equals $E$ near $C_U$, so it glues to an actual form on $U$.
The relative-support Stokes calculation of the companion proves
the all-actual test identity.

For (ii), lift the fixed finite word locally and execute its
factors at the actual intermediate tuples. The final discrepancy
is $1+E_0$ with entries of $E_0$ vanishing on the reduced common
zero scheme. Near $C_U$ its entries have norm less than one, so
this discrepancy preserves the tuple supremum norm and the
complete expression \eqref{eq:ci-kernel}. The coefficient-uniform
germ theorem in \cite{InternalTropical} identifies the complete
primitive germs for different local lifts. They therefore glue
with a locally finite partition; its derivative discrepancy
extends by zero. Apply the same relative-support cutoff argument.
The calculation retains the original nonreduced schemes and their
cycle multiplicities throughout.
\end{proof}

These regular sections have
$[Y_i]=c_r(L^{\oplus r})\cap[X]=c_1(L)^r\cap[X]$ in Chow.
Consequently $[Z]=0$. The constructions make the exact-support
Green representative explicit. For such globally defined cycles,
the Chow-descended height is
\[
 h_X([Z],w)=0\qquad(w\in\cH^{n+1-r}(X)),
\]
even though individual local Green integrals may be nonzero.

\subsection{Comparison with supplied arithmetic models}

Suppose an admitted pair $z,w$ has regular numerically admissible
arithmetic lifts $\widehat\Theta,\widehat\Phi$ on a regular flat
projective model over $\mathcal O_F$, with zero archimedean
curvature in the comparison normalization. The companion
model-comparison theorem \cite{InternalHeight} identifies its
normalized height with
\[
 h_X(z,w)=[F:\Q]^{-1}
          \widehat{\deg}(\widehat\Theta\cdot\widehat\Phi).
\]
If $f:Y\to X_E$ is a degree-$d$ alteration with $Y$ smooth
projective, and the supplied admissible lifts of
$f^*z_E,f^*w_E$ lie on a regular flat projective
$\mathcal O_E$-model of $Y$, with the same zero-curvature
archimedean normalization, the corresponding formula is
\[
 h_X(z,w)=\frac{1}{d[E:\Q]}
          \widehat{\deg}(\widehat\Theta\cdot\widehat\Phi).
\]
The generic degree and the absolute number-field degree are separate
factors. These formulas and the abelian comparison in the next
section describe the same normalized pairing on their common domains.
